\section{Introduction}

The present paper is not anti-G\"odel. Classical incompleteness remains intact where it ordinarily lives: in standard metamathematical settings with arithmetized syntax, provability predicates, and diagonal construction. The question here is narrower. Suppose one already works inside a fixed same-domain regime forced into a hard kernel of inferential governance: standing and admissibility are bivalent, the admissible interior is unique, lawful redescription is tensor-preserving, coherent trajectories are pointwise standing-preserving legal paths, no deferred coherence is possible, terminal standing fully certifies coherence, and no same-scope repair, generator, or independent authorizer survives within the regime. Can a G\"odel-style threat architecture still be instantiated from within that same domain as a governance-relevant closure threat?

This paper answers no.

The proof strategy is eliminative rather than revisionary. It does not deny the classical incompleteness theorems. It asks instead what a same-domain incompleteness objection would have to do in order to matter \emph{here}. Such an objection would have to change inferential life on the original fixed domain: what may be admitted, rejected, reused, committed, or licensed as same-domain continuation. The first bridge theorem therefore shows that inferential life is exhausted by two and only two loci: standing classification of evaluated acts, and coherent legal trajectory structure of same-domain histories. The coherent reasoning layer is decisive here. Once coherence is pointwise standing-preservation, once incoherent prefixes remain incoherent under admissible continuation, and once terminal standing fully certifies trajectory coherence, there is no third hidden path-locus in which a same-domain closure threat can matter. This is also why the target is not bare derivation-trace classification. A syntax, coding, or proof inscription matters in the present paper only if its same-domain readback changes standing classification or coherent legal trajectory structure. If it does neither, it is not a closure threat but bookkeeping about a closure threat.

The body then applies that bridge mechanically. Governance-relevant code predicates become hidden gate work. Same-domain negative fixed-point routes fail because no closure-relevant distinction survives their readback. Semantic truth, reflection, infinitary rule-power, meta-foundational re-anchoring, and richer same-domain extension do not rescue the architecture: they either collapse to bookkeeping, import illicit same-domain authority, or become explicit scope change. The terminal theorem therefore states not that G\"odel is false, but that a G\"odel-style \emph{same-domain} threat architecture is not instantiable inside the unique admissible interior once the coherent-reasoning kernel is in force.


A full technical appendix package now supports the primary manuscript directly, while the fully armored audit edition is retained separately as an audit companion rather than being allowed to dominate the visible body. Formalization is treated in the same way: the prose theorem chain is primary, and detailed theorem-ledger provenance is relegated to the audit companion.